\documentclass[10pt]{extarticle}
\usepackage[a4paper,margin=1.5cm]{geometry}
\usepackage[T1]{fontenc}
\usepackage{amsmath,amssymb}
\usepackage{graphicx}
\usepackage[export]{adjustbox}
\usepackage{multicol}
\usepackage{enumitem}
\usepackage{booktabs,array,tabularx}
\setlength{\columnsep}{0.7cm}
\setlength{\parindent}{0pt}
\setlength{\parskip}{3pt}
\pagestyle{empty}
\makeatletter
% figure registry (figdefs.tex is generated by the extract step)
\newcommand{\qfigdef}[4]{\@namedef{qf@file@#1}{#2}\@namedef{qf@w@#1}{#3}\@namedef{qf@wide@#1}{#4}}
\newcommand{\qf@check}[1]{\@ifundefined{qf@file@#1}{\PackageError{score2}{Unknown figure id #1}{}}{}}
\newcommand{\qf@img}[4][]{\includegraphics[width=#2\dimexpr\@nameuse{qf@w@#4}pt\relax,max width=#3,max height=0.45\textheight,#1]{\@nameuse{qf@file@#4}}}
% Full-width questions: finalize wraps every question that contains a wide figure in \qfullbegin...\qfullend,
% so the WHOLE question (stem, figures, options) is set across the page in reading order, then the two columns
% resume. The question's head - its number, stem and first figure (\qfullsplit marks the end of the head) -
% is kept on one page, so a number is never left at a page foot with its figure overleaf; the rest may break
% across pages normally, so no big gaps are left. Without a \qfullsplit, the whole question is kept together
% if it is under 45% of a page. Never write these yourself.
\newif\ifqf@full
\newif\ifqf@boxed
\newsavebox{\qf@box}
\newcommand{\qf@open}{\global\qf@fulltrue\par\global\qf@boxedtrue
  \setbox\qf@box\vbox\bgroup\hsize\textwidth\linewidth\textwidth}
\newcommand{\qf@emit}[1]{\par\egroup\global\qf@boxedfalse
  \ifdim\dimexpr\ht\qf@box+\dp\qf@box\relax<#1\textheight\noindent\usebox{\qf@box}\par\else\unvbox\qf@box\fi}
\newcommand{\qfullsplit}{\ifqf@boxed\qf@emit{0.85}\fi}
\newcommand{\qf@close}{\ifqf@boxed\qf@emit{0.45}\fi\global\qf@fullfalse}
\newcommand{\qfullbegin}{\end{multicols}\qf@open}
\newcommand{\qfullend}{\qf@close\begin{multicols}{2}}
\newcommand{\qfullstart}{\qf@open}
\newcommand{\qfulldone}{\qf@close}
% \qfig[scale]{id}: block figure, centred, directly where it is written (size set by the tool).
% Fallback if a wide figure is not inside a full-width question: break out of the columns in place - no floats.
\newcommand{\qfig}[2][1]{\qf@check{#2}%
  \ifqf@full\qf@blk{#1}{#2}\else
  \ifnum\@nameuse{qf@wide@#2}=1 \expandafter\@firstoftwo\else\expandafter\@secondoftwo\fi
  {\end{multicols}\begin{center}\qf@img{#1}{\textwidth}{#2}\end{center}\begin{multicols}{2}}%
  {\qf@blk{#1}{#2}}\fi}
\newcommand{\qf@blk}[2]{\par\smallskip{\centering\qf@img{#1}{\linewidth}{#2}\par}\smallskip}
% \qfigi[scale]{id}: inline figure (inside option lists / tables / minipages); never breaks columns.
\newcommand{\qfigi}[2][1]{\qf@check{#2}\qf@img[valign=c]{#1}{\linewidth}{#2}}
\makeatother
% structure macros - lint reads these to check completeness and figure placement
\newcommand{\papersec}[1]{\par\bigskip{\centering\Large\bfseries PAPER-#1\par}\medskip}
\newcommand{\partsec}[2]{\par\medskip\noindent{\large\bfseries PART-#1: #2}\par\smallskip}
\newcommand{\sect}[1]{\par\smallskip\noindent\textbf{#1}\par}
\newcommand{\qn}[1]{\par\smallskip\noindent\textbf{#1.}~}
\newcommand{\paraq}[2]{\par\smallskip\noindent\textbf{Paragraph for Questions #1 and #2}\par}
\newenvironment{qoptions}{\begin{itemize}[leftmargin=1.8em,itemsep=1pt,topsep=2pt,label={}]}{\end{itemize}}
\newcommand{\opt}[2]{\item[(#1)] #2}
% figure registry (generated by score2_tools.py finalize)
\qfigdef{p1_phy_s1_q01_1}{Star Review (Other) - Test 01 (12 Sep 2022)-img/p1_phy_s1_q01_1.png}{121.7}{0}
\qfigdef{p1_phy_s1_q04_1}{Star Review (Other) - Test 01 (12 Sep 2022)-img/p1_phy_s1_q04_1.png}{163.9}{0}
\qfigdef{p1_phy_s1_q05_1}{Star Review (Other) - Test 01 (12 Sep 2022)-img/p1_phy_s1_q05_1.png}{81.6}{0}
\qfigdef{p1_phy_s1_q06_1}{Star Review (Other) - Test 01 (12 Sep 2022)-img/p1_phy_s1_q06_1.png}{124.2}{0}
\qfigdef{p1_phy_s2_para7_8_1}{Star Review (Other) - Test 01 (12 Sep 2022)-img/p1_phy_s2_para7_8_1.png}{142.8}{0}
\qfigdef{p1_phy_s2_para9_10_1}{Star Review (Other) - Test 01 (12 Sep 2022)-img/p1_phy_s2_para9_10_1.png}{134.9}{0}
\qfigdef{p1_phy_s3_para1_2_1}{Star Review (Other) - Test 01 (12 Sep 2022)-img/p1_phy_s3_para1_2_1.png}{92.6}{0}
\qfigdef{p1_phy_s3_para3_4_1}{Star Review (Other) - Test 01 (12 Sep 2022)-img/p1_phy_s3_para3_4_1.png}{116.5}{0}
\qfigdef{p1_phy_s3_para5_6_1}{Star Review (Other) - Test 01 (12 Sep 2022)-img/p1_phy_s3_para5_6_1.png}{95.0}{0}
\qfigdef{p1_phy_s4_q07_1}{Star Review (Other) - Test 01 (12 Sep 2022)-img/p1_phy_s4_q07_1.png}{137.6}{0}
\qfigdef{p1_phy_s4_q09_1}{Star Review (Other) - Test 01 (12 Sep 2022)-img/p1_phy_s4_q09_1.png}{145.2}{0}
\qfigdef{p1_chem_s1_q05_A}{Star Review (Other) - Test 01 (12 Sep 2022)-img/p1_chem_s1_q05_A.png}{63.5}{0}
\qfigdef{p1_chem_s1_q05_B}{Star Review (Other) - Test 01 (12 Sep 2022)-img/p1_chem_s1_q05_B.png}{82.1}{0}
\qfigdef{p1_chem_s1_q05_C}{Star Review (Other) - Test 01 (12 Sep 2022)-img/p1_chem_s1_q05_C.png}{58.3}{0}
\qfigdef{p1_chem_s2_para7_8_1}{Star Review (Other) - Test 01 (12 Sep 2022)-img/p1_chem_s2_para7_8_1.png}{34.6}{0}
\qfigdef{p1_chem_s2_para7_8_2}{Star Review (Other) - Test 01 (12 Sep 2022)-img/p1_chem_s2_para7_8_2.png}{58.3}{0}
\qfigdef{p1_chem_s3_para1_2_1}{Star Review (Other) - Test 01 (12 Sep 2022)-img/p1_chem_s3_para1_2_1.png}{121.0}{0}
\qfigdef{p1_chem_s3_para1_2_2}{Star Review (Other) - Test 01 (12 Sep 2022)-img/p1_chem_s3_para1_2_2.png}{97.0}{0}
\qfigdef{p2_phy_s1_q02_1}{Star Review (Other) - Test 01 (12 Sep 2022)-img/p2_phy_s1_q02_1.png}{124.2}{0}
\qfigdef{p2_phy_s1_q03_1}{Star Review (Other) - Test 01 (12 Sep 2022)-img/p2_phy_s1_q03_1.png}{131.5}{0}
\qfigdef{p2_phy_s1_q04_1}{Star Review (Other) - Test 01 (12 Sep 2022)-img/p2_phy_s1_q04_1.png}{131.3}{0}
\qfigdef{p2_phy_s1_q05_1}{Star Review (Other) - Test 01 (12 Sep 2022)-img/p2_phy_s1_q05_1.png}{120.8}{0}
\qfigdef{p2_phy_s1_q06_1}{Star Review (Other) - Test 01 (12 Sep 2022)-img/p2_phy_s1_q06_1.png}{97.9}{0}
\qfigdef{p2_phy_s1_q07_1}{Star Review (Other) - Test 01 (12 Sep 2022)-img/p2_phy_s1_q07_1.png}{79.4}{0}
\qfigdef{p2_phy_s1_q09_1}{Star Review (Other) - Test 01 (12 Sep 2022)-img/p2_phy_s1_q09_1.png}{76.8}{0}
\qfigdef{p2_phy_s1_q15_1}{Star Review (Other) - Test 01 (12 Sep 2022)-img/p2_phy_s1_q15_1.png}{111.2}{0}
\qfigdef{p2_phy_s1_q16_1}{Star Review (Other) - Test 01 (12 Sep 2022)-img/p2_phy_s1_q16_1.png}{147.5}{0}
\qfigdef{p2_phy_s1_q17_1}{Star Review (Other) - Test 01 (12 Sep 2022)-img/p2_phy_s1_q17_1.png}{63.5}{0}
\qfigdef{p2_phy_s1_q19_1}{Star Review (Other) - Test 01 (12 Sep 2022)-img/p2_phy_s1_q19_1.png}{92.3}{0}
\qfigdef{p2_phy_s2_q01_1}{Star Review (Other) - Test 01 (12 Sep 2022)-img/p2_phy_s2_q01_1.png}{144.9}{0}
\qfigdef{p2_phy_s2_q04_1}{Star Review (Other) - Test 01 (12 Sep 2022)-img/p2_phy_s2_q04_1.png}{147.5}{0}
\qfigdef{p2_phy_s2_q05_1}{Star Review (Other) - Test 01 (12 Sep 2022)-img/p2_phy_s2_q05_1.png}{150.1}{0}
\qfigdef{p2_phy_s2_q08_1}{Star Review (Other) - Test 01 (12 Sep 2022)-img/p2_phy_s2_q08_1.png}{152.6}{0}
\qfigdef{p2_phy_s2_q09_1}{Star Review (Other) - Test 01 (12 Sep 2022)-img/p2_phy_s2_q09_1.png}{153.9}{0}
\qfigdef{p2_phy_s2_q10_1}{Star Review (Other) - Test 01 (12 Sep 2022)-img/p2_phy_s2_q10_1.png}{121.3}{0}
\qfigdef{p2_chem_s2_q05_1}{Star Review (Other) - Test 01 (12 Sep 2022)-img/p2_chem_s2_q05_1.png}{31.8}{0}
\qfigdef{p2_chem_s2_q05_2}{Star Review (Other) - Test 01 (12 Sep 2022)-img/p2_chem_s2_q05_2.png}{21.3}{0}
\qfigdef{p2_chem_s2_q09_1}{Star Review (Other) - Test 01 (12 Sep 2022)-img/p2_chem_s2_q09_1.png}{29.2}{0}
\raggedcolumns
\makeatletter\@beginparpenalty=10000\makeatother
\newcolumntype{L}[1]{>{\raggedright\arraybackslash}p{#1}}
\begin{document}
\begin{center}{\Large\bfseries Star Review (Other) --- Test 01 (12 Sep 2022)}\end{center}
\begin{multicols}{2}
\raggedcolumns
\papersec{1}
\partsec{1}{PHYSICS}
\sect{SECTION-I(i)}
\qn{1} Two identical uniform equilateral triangular plates are joined as shown. If $(x,y)$ are the coordinates of centre of mass of this combined plate, then :
\qfig{p1_phy_s1_q01_1}
\begin{qoptions}
\opt{A}{$x=\dfrac{\ell\sqrt3}2$}
\opt{B}{$x=\dfrac\ell2$}
\opt{C}{$y=\dfrac\ell2$}
\opt{D}{$y=\dfrac{\ell\sqrt3}2$}
\end{qoptions}
\qn{2} A body of mass $m=1$kg is displaced by force $\vec F=\left(2y\hat i+3x^2\hat j\right)$ N from $\vec r_1=\left(\hat i+\hat j\right)$ m to $\vec r_2=\left(2\hat i+4\hat j\right)$ m along the path $y=x^2$. If initially body was at rest, then
\begin{qoptions}
\opt{A}{work-done by the force is $\dfrac{163}6$J}
\opt{B}{work-done by the force is $\dfrac{163}3$J.}
\opt{C}{Final speed of the body is $\sqrt{\dfrac{163}3}$m/s.}
\opt{D}{Final speed of the body is $\sqrt{\dfrac{163}6}$m/s.}
\end{qoptions}
\qn{3} If force $(F)$, momentum $(P)$ and mass $(M)$ are fundamental quantities in system of units, then dimensional formula for length $(L)$ and time $(T)$ will be :-
\begin{qoptions}
\opt{A}{$[T]=[F^2P^{-1}M^0]$}
\opt{B}{$[L]=[F^{-1}M^{-2}P^{-1}]$}
\opt{C}{$[L]=[F^{-1}P^2M^{-1}]$}
\opt{D}{$[T]=[F^{-1}P^1M^0]$}
\end{qoptions}
\qn{4} A small block of mass 3 kg is projected into a vertical return loop at A with velocity $v$. The block travels without friction along a circle of radius 2m, and deposited on a horizontal surface at C for each of two loops shown :
\qfig{p1_phy_s1_q04_1}
\begin{qoptions}
\opt{A}{In case-I, the smallest velocity for which the block will reach the horizontal surface at C is given by 10 m/s.}
\opt{B}{In case-II, the smallest velocity for which the block will reach the horizontal surface at C is given by $\sqrt{80}$m/s.}
\opt{C}{In case-I force exerted by the loop on the block as it passes point B is given by 90 N.}
\opt{D}{In case-II force exerted by the loop on the block as it passes point B is given by 60 N.}
\end{qoptions}
\qn{5} Two blocks of masses $m_1=2$ kg and $m_2=4$ kg hang over a massless pulley as shown in the figure. A force $F_0=100$ N acting at the axis of the pulley accelerates the system upwards. Then :
\qfig{p1_phy_s1_q05_1}
\begin{qoptions}
\opt{A}{Acceleration of 2 kg mass is 15 m/s$^2$}
\opt{B}{Acceleration of 4 kg mass is 2.5 m/s$^2$}
\opt{C}{Acceleration of both the masses is same}
\opt{D}{Acceleration of both the masses is upward.}
\end{qoptions}
\qn{6} One end of a thread of length `$\ell$' is attached to an inclined plane of elevation of $\alpha$. A point like body of mass $m=1$ kg is attached to the other end of the thread as shown in the figure. The body is released without initial speed, such that the thread is tight and horizontal. The coefficient of friction between the slope and body is $\mu$.
\qfig{p1_phy_s1_q06_1}
\begin{qoptions}
\opt{A}{The tension in the thread, when the angle between the thread and the horizontal is $\phi$, is $T=mg\left(3\sin\phi\sin\alpha-2\mu\phi\cos\alpha\right)$.}
\opt{B}{The tension in the string will be maximum when $\phi=90^\circ$.}
\opt{C}{The tension in the string will be maximum when $\phi=\cos^{-1}\left(\dfrac{2\mu}{3\tan\alpha}\right)$.}
\opt{D}{The maximum speed of the particle will be, when $\phi=\cos^{-1}(\mu\cot\alpha)$}
\end{qoptions}
\sect{SECTION-I(ii)}
\paraq{7}{8}
A small block of mass $M$, starts sliding down on a bigger block of mass $2M$, having a circular track of radius $R$. Let $N_1$ and $N_2$ be the normal reactions on the bigger block due to the wall and the floor. Do not consider the motion of the small block after $\theta=90^\circ$. (All surfaces are frictionless)
\qfig{p1_phy_s2_para7_8_1}
\qn{7} The value of $N_1$ at the instant when $\theta=30^\circ$
\begin{qoptions}
\opt{A}{$\dfrac{\sqrt3Mg}4$}
\opt{B}{$\dfrac{3\sqrt3Mg}4$}
\opt{C}{$\dfrac{\sqrt3Mg}2$}
\opt{D}{$\dfrac{3\sqrt3Mg}2$}
\end{qoptions}
\qn{8} The value of $\tan\theta$ at the instant when $N_2=3\,Mg$
\begin{qoptions}
\opt{A}{$\dfrac1{\sqrt2}$}
\opt{B}{$\sqrt2$}
\opt{C}{$\dfrac12$}
\opt{D}{$2$}
\end{qoptions}
\paraq{9}{10}
A ball of mass $M$ attached to a string of length $\ell$ moves in a circle in a vertical plane as shown. At the top of the circular path, the tension is twice the weight of the ball. At the bottom, the ball just clear the ground.
\qfig{p1_phy_s2_para9_10_1}
\qn{9} Determine the speed of the ball at the lower most point.
\begin{qoptions}
\opt{A}{$\sqrt{3g\ell}$}
\opt{B}{$\sqrt{5g\ell}$}
\opt{C}{$\sqrt{6g\ell}$}
\opt{D}{$\sqrt{7g\ell}$}
\end{qoptions}
\qn{10} If the string is cut when the ball is at the top at $t=t_0$, then the horizontal distance the ball travels before hitting the ground and after $t=t_0$ is :-
\begin{qoptions}
\opt{A}{$2\sqrt3\ell$}
\opt{B}{$4\sqrt3\ell$}
\opt{C}{$3\sqrt3\ell$}
\opt{D}{$5\sqrt3\ell$}
\end{qoptions}
\sect{SECTION-II(i)}
\paraq{1}{2}
A man of mass 50 kg stands on a horizontal weighing machine of negligible mass, attached to massless platform that slide down a $37^\circ$ rough incline plane. The reading of weighing machine is 48 kg (take : $g=10$ m/s$^2$)
\qfig{p1_phy_s3_para1_2_1}
\qn{1} What is the vertical acceleration (in m/s$^2$) of man ?
\qn{2} What is the acceleration (in m/s$^2$) of man down the slope $a_s$ ?
\paraq{3}{4}
Masses of blocks A, B and C are 7kg, 4kg and 1 kg respectively. String and pulleys are ideal. All surfaces are smooth.
\qfig{p1_phy_s3_para3_4_1}
\qn{3} Find the tension in the string (in N).
\qn{4} The acceleration of block `C' is $\sqrt N$ m/s$^2$. Find `N'.
\paraq{5}{6}
A sufficiently long plank of mass 4 kg is kept on a smooth horizontal ground. A block of mass 2kg kept over the plank and is being acted upto by a time varying horizontal force $F=0.5\,t$, where $F$ is in newton and $t$ is in seconds as shown in figure. The coefficient of friction between the plank and the block is `$\mu$'. At $t=12$ sec relative motion between the plank and the block just starts.
\qfig{p1_phy_s3_para5_6_1}
\qn{5} Find the coefficient of friction $\mu$.
\qn{6} Find the average acceleration of the plank in the time interval 0 to 15 s in the figure (in m/s$^2$).
\sect{SECTION-II(ii)}
\qn{7} A particle moves in rectilinear motion. The acceleration-position curve is shown below. If initial velocity of particle is 20 m/s. Velocity of particle (in m/s) after the particle has travelled 50 m is
\qfig{p1_phy_s4_q07_1}
\qn{8} Velocity of river flow grows in proportion to the distance from the bank and reaches its maximum velocity of 2 m/s in the middle. Near the bank, velocity is zero. The velocity of a swimmer in still water is 5 m/s and is directed perpendicular to river flow. The width of the river is 100 m. Find the drift (in m) of the swimmer.
\qn{9} As shown in the arrangement a uniform heavy rope AB of mass 10 kg slides down smooth surface. Find the tension (in N) at $\dfrac34$th length of rope from point A.
\qfig{p1_phy_s4_q09_1}
\partsec{2}{CHEMISTRY}
\sect{SECTION-I(i)}
\qn{1} 5 moles of $\mathrm{N_2}$ and 20 moles of $\mathrm{H_2}$ are allowed to react to form $\mathrm{NH_3}$ , as the only product. The correct information(s) is/are
\begin{qoptions}
\opt{A}{At $t=0$, average molar mass of gaseous mixture $=7.2$ gm/mol}
\opt{B}{After complete reaction, average molar mass of gaseous mixture $=12$ gm /mol}
\opt{C}{After complete reaction 5gm molecule of $\mathrm{H_2}$ left unreacted}
\opt{D}{At any stage of reaction, vapour density of gas mixture is constant}
\end{qoptions}
\qn{2} For 2.6 kg of $\mathrm{FeSO_4\cdot6H_2O}$, the correct information(s) is/are (At. wt. of Fe = 56)
\begin{qoptions}
\opt{A}{Number of atoms of O are $100\times N_A$}
\opt{B}{Moles of H-atoms are $120\times N_A$}
\opt{C}{Number of molecules of water of crystalline are $60\times N_A$}
\opt{D}{Moles of electron present in $\mathrm{SO_4^{2-}}$ ions in sample are 500}
\end{qoptions}
\qn{3} 12g of Mg was burnt in a closed vessel containing 32g oxygen. Which of the following is/are correct.
\begin{qoptions}
\opt{A}{2gm of Mg will be left unburnt}
\opt{B}{0.75 gm-molecule of $\mathrm{O_2}$ will be left unreacted}
\opt{C}{20gm of MgO will be formed}
\opt{D}{The mixture at the end will weight 44g}
\end{qoptions}
\qn{4} Which of the following set(s) have I$^{\mathrm{st}}$ element with both EA \& IE higher than 2$^{\mathrm{nd}}$ element.
\begin{qoptions}
\opt{A}{(N, O)}
\opt{B}{(P, Si)}
\opt{C}{(F, O)}
\opt{D}{(S, Se)}
\end{qoptions}
\qn{5} Which of the following interactions have energy proportional to $\dfrac1{r^3}$ (`$r$' is separation between species)?
\begin{qoptions}
\opt{A}{\qfigi{p1_chem_s1_q05_A}}
\opt{B}{\qfigi{p1_chem_s1_q05_B}}
\opt{C}{\qfigi{p1_chem_s1_q05_C}}
\opt{D}{$\mathrm{CH_3Cl+CCl_4}$}
\end{qoptions}
\qn{6} Among the following cases, the ion(s)/molecules that have X--O--X linkage ($X\to$ so called central atom)
\begin{qoptions}
\opt{A}{$\mathrm{P_4O_{10}}$}
\opt{B}{$\mathrm{H_2P_2O_7^{2-}}$}
\opt{C}{$\mathrm{C_2O_4^{2-}}$}
\opt{D}{$\mathrm{Cr_2O_7^{2-}}$}
\end{qoptions}
\sect{SECTION-I(ii)}
\paraq{7}{8}
Double chain structures of $\{\mathrm{SiO_4}\}$ tetrahedra are given below.
\qfig{p1_chem_s2_para7_8_1}
\qfig{p1_chem_s2_para7_8_2}
\qn{7} Which of the following correctly depicts the general formula unit for I
\begin{qoptions}
\opt{A}{$\mathrm{Si_2O_4^{2-}}$}
\opt{B}{$\mathrm{Si_2O_3^{2-}}$}
\opt{C}{$\mathrm{Si_2O_5^{2-}}$}
\opt{D}{$\mathrm{SiO_5^{2-}}$}
\end{qoptions}
\qn{8} Among the given species, which one has more number of atoms having tetrahedral molecular geometry around them as compared to smallest repeating unit of II
\begin{qoptions}
\opt{A}{$\mathrm{C(CN)_4}$}
\opt{B}{$\mathrm{C_2H_2(CH_3)_2}$}
\opt{C}{$\mathrm{H_4P_2O_7}$}
\opt{D}{$\mathrm{C(CH_3)_4}$}
\end{qoptions}
\paraq{9}{10}
An element "X" of 3d series has maximuim number of unpaired electrons among all elements in its period. The element "Y" adjacent to it forms a monobasic oxyacid having tetrahedral shape.
\qn{9} Which is correct regarding both elements.
\begin{qoptions}
\opt{A}{$\mathrm{X_2O_3}$ is acidic}
\opt{B}{$\mathrm{X^{3+}}$ is diamagnetic}
\opt{C}{$\mathrm{Y_2O_7}$ is amphoteric}
\opt{D}{$\mathrm{Y^{2+}}$ is paramagnetic}
\end{qoptions}
\qn{10} Which of the following statements is INCORRECT?
\begin{qoptions}
\opt{A}{$\mathrm{XO_4^{2-}}$ \& $\mathrm{YO_4^{2-}}$ are isostructural.}
\opt{B}{X has more number of possible exchange pairs compared to Y in "d" subshell.}
\opt{C}{$\mathrm{X^+}$ \& $\mathrm{Y^{2+}}$ are isoelectronic}
\opt{D}{Both X and Y can form basic, amphoteric \& acidic oxides}
\end{qoptions}
\sect{SECTION-II(i)}
\paraq{1}{2}
Born - Haber cycle for formation of solid binary metal hydride is given below along with energy profile diagram of $\mathrm{H_2}$ molecule.
\qfig{p1_chem_s3_para1_2_1}
where Me $\to$ Alkali metal or Alkaline earth metal
\qfig{p1_chem_s3_para1_2_2}
The following data can be useful.\par
$\left(\Delta H_{IE_1}\right)_{Na}=495$ kJ/mol\par
$\left(\Delta H_{IE_1}\right)_{Ca}=589$ kJ/mol\par
$\left(\Delta H_{EG}\right)_{Ca^{2+}}=-1145$ kJ/mol\par
$\left(\Delta H_{LE}\right)_{NaH}=808$ kJ/mol\par
$\left(\Delta H_{LE}\right)_{CaH_2}=2390$ kJ/mol\par
$\left(\Delta H_{E.G._1}\right)_H=-72$ kJ/mol
\qn{1} The value of $\left|\left(\Delta H_f\right)_{CaH_2}-\left(\Delta H_f\right)_{NaH}\right|$ if $\left\{\left(\Delta H_{sublimation}\right)_{Ca}-\left(\Delta H_{Sublimation}\right)_{Na}\right\}$ is 200 kJ/mol is \_\_\_\_\_\_.
\qn{2} If change in enthalpy for the process $\mathrm{Me(g)}+\dfrac n2\mathrm H_2(g)\to\mathrm{Me^{n+}(g)}+n\mathrm H^-(g)$ is $\Delta H_1$ then the value of $|\Delta H_1|$ when Me is "Na" is \_\_\_\_\_\_\_
\paraq{3}{4}
An element "A" forms a compound I comprising of A, F and O having see-saw shape. Another compound II of A, F \& O having non-zero dipole moment has adjacent $\mathrm F-\hat A-\mathrm F$ bond angle significantly lower than $\mathrm F-\hat A-\mathrm F$ bond angle in see-saw shape compound.
\qn{3} If "A" is Xe, then difference in number of lone pairs in both compounds is \_\_\_\_\_\_\_\_\_\_\_\_\_.
\qn{4} If compound II has zero dipole moment then difference in number of lone pairs in both compounds is \_\_\_\_\_\_\_\_.
\paraq{5}{6}
Air is a mixture of 80\% $\mathrm{N_2}$ and 20\% $\mathrm{O_2}$ by moles. One mole of mixture of air and $\mathrm{CH_4}$ is taken and subjected to combustion of $\mathrm{CH_4}$. It is also given that 1 mole of $\mathrm{CH_4}$ on complete combustion gives 220 kJ of heat.
\qn{5} $x$, $y$, $z$ are the moles of each component $\mathrm{O_2}$, $\mathrm{N_2}$ and $\mathrm{CH_4}$ so as to release maximum heat then find value of $\left(\dfrac{11xy}z\right)$.
\qn{6} What will be the maximum heat released in kJ ?
\sect{SECTION-II(ii)}
\qn{7} 50\% pure sample of $\mathrm{CaCO_3}$ heated with 80\% yield to produce 0.448 litre $\mathrm{CO_2}$ at 1 atm and 273 K temperature. Calculate mass of $\mathrm{CaCO_3}$ sample in gm.
\qn{8} The value of `x', if hydrated salt $\mathrm{A_2SO_4\cdot xH_2O}$ undergoes 45\% loss in mass on heating and becomes anhydrous, is (where atomic weight of `A' is 7)
\qn{9} Consider three neutral molecules of the type $\mathrm{AB_2L_2}$ , $\mathrm{ABL_3}$ and $\mathrm{AB_3L}$.\par
($L\to$ Lone pair on A)\par
Among $\mathrm{XeO_3}$ , $\mathrm{SO_2}$ , $\mathrm{XeF_2}$ , $\mathrm{ClF_3}$ , $\mathrm{SnCl_3^{\ominus}}$ , $\mathrm{SO_3}$ , $\mathrm{ICl}$ , $\mathrm{H_2O}$\par
$X\to$ number of $\mathrm{AB_2L_2}$ type molecules\par
$Y\to$ number of $\mathrm{ABL_3}$ type molecules\par
$Z\to$ number of $\mathrm{AB_3L}$ type molecules\par
then $(X+Y+Z)$ is \_\_\_\_\_\_\_\_\_\_.
\partsec{3}{MATHEMATICS}
\sect{SECTION-I(i)}
\qn{1} If roots of the equation $x^3+\alpha x^2+3x+4=0$ are $\sec\theta_1+\tan\theta_1$, $\sec\theta_2+\tan\theta_2$, $\sec\theta_3+\tan\theta_3$ and equation whose roots are $\sec\theta_1-\tan\theta_1$, $\sec\theta_2-\tan\theta_2$, $\sec\theta_3-\tan\theta_3$ are $x^3+\beta x^2+\gamma x+\delta=0$, then which of the following option(s) are divisor of $4\beta+8\delta$ is
\begin{qoptions}
\opt{A}{$0$}
\opt{B}{$3$}
\opt{C}{$5$}
\opt{D}{$15$}
\end{qoptions}
\qn{2} If $\sin x=\dfrac{\left(\sqrt5-1\right)}4$, then which of the following is/are correct ?
\begin{qoptions}
\opt{A}{$\cos4x=\sin x$}
\opt{B}{$\sin2x+\cos3x=\dfrac12$}
\opt{C}{$\sin3x=\cos2x$}
\opt{D}{$\sin5x+\cos5x=1$}
\end{qoptions}
\qn{3} Which of the following interval(s) will be the subset of the solution set of the inequality $\left(x^2-x-1\right)\left(x^2-x-7\right)<-5$
\begin{qoptions}
\opt{A}{$(-1,3)$}
\opt{B}{$(-2,-1)$}
\opt{C}{$(-2,3)$}
\opt{D}{$(2,3)$}
\end{qoptions}
\qn{4} In $\Delta ABC$, if $\dfrac1{a+c}+\dfrac1{b+c}=\dfrac3{a+b+c}$ then which of the following hold good ?
\begin{qoptions}
\opt{A}{$\displaystyle\sum_{r=1}^{\infty}\cos^{2r}(A+B)=1$}
\opt{B}{$\displaystyle\sum_{r=1}^{\infty}\cos^{2r}(A+B)=\dfrac13$}
\opt{C}{$\dfrac{\cot C}{\cot A+\cot B}$ is maximum if $\Delta ABC$ is equilateral}
\opt{D}{$\dfrac{\cot C}{\cot A+\cot B}$ is maximum if $\Delta ABC$ is isosceles}
\end{qoptions}
\qn{5} If $P=\left\{\theta;\ \sin\theta<\dfrac{-1}{\sqrt2}\right\}$\par
$Q=\left\{\theta;\ \cos\theta<\dfrac{-1}2\right\}$\par
$R=\{\theta;\ \tan\theta>1\}$\par
then , which of the following is/are hold good
\begin{qoptions}
\opt{A}{$P\cap Q\cap R\in\left(\dfrac{5\pi}4,\dfrac{4\pi}3\right)$}
\opt{B}{$P\cap Q\cap R\in\left(\dfrac{-3\pi}4,\dfrac{-2\pi}3\right)$}
\opt{C}{$P\cap Q\in\left(\dfrac{5\pi}4,\dfrac{4\pi}3\right)$}
\opt{D}{$Q\cap R\in\left(\dfrac{-3\pi}4,\dfrac{-2\pi}3\right)$}
\end{qoptions}
\qn{6} If the value of expression $\sin\left(\dfrac{3\pi}{30}\right)$ $\sin\left(\dfrac{7\pi}{30}\right)$ $\sin\left(\dfrac{13\pi}{30}\right)$ can be expressed as $\dfrac{\sqrt a+\sqrt b}c$, where $a,b,c\in N$, $\{a>b\}$ and are in their lowest form, then
\begin{qoptions}
\opt{A}{$c-a-b=10$}
\opt{B}{$a+b+c=21$}
\opt{C}{$c\div b-a=11$}
\opt{D}{$c\div b-a=-4$}
\end{qoptions}
\sect{SECTION-I(ii)}
\paraq{7}{8}
Consider the series\par
{\scriptsize$\underbrace{1^2x^1}_{\text{1 term}}+\underbrace{2^2x^2+2^2x^3}_{\text{2 terms}}+$}\par {\scriptsize$\underbrace{4^2x^4+4^2x^5+4^2x^6+4^2x^7}_{\text{4 terms}}+$}\par {\scriptsize$\underbrace{8^2x^8+8^2x^9+\dots+\dots}_{\text{8 terms}}+\dots$}\par
Now answer the following
\qn{7} If $x=1$, then the sum of first 1023 terms is equal to
\begin{qoptions}
\opt{A}{$\dfrac{4^{20}-1}3$}
\opt{B}{$\dfrac{8^{10}-1}7$}
\opt{C}{$\left(2^{1023}-1\right)$}
\opt{D}{$16^5$}
\end{qoptions}
\qn{8} If $x=-1$, then the sum of first 1023 terms is equal to
\begin{qoptions}
\opt{A}{$-1$}
\opt{B}{$\dfrac{2^{20}-1}3$}
\opt{C}{$\dfrac{4^{10}-1}3$}
\opt{D}{$\dfrac{8^{10}-1}7$}
\end{qoptions}
\paraq{9}{10}
If $P=\{x;\ \cos3x=\cos x\ \&\ x\in[0,2\pi]\}$ \& $Q=\{x;\ \tan4x=\tan x\ \&\ x\in[0,2\pi]\}$ then
\qn{9} $n(P\cap Q)$ is \_\_\_\_\_\_\_\_\_
\begin{qoptions}
\opt{A}{$1$}
\opt{B}{$2$}
\opt{C}{$3$}
\opt{D}{$4$}
\end{qoptions}
\qn{10} $n(P\cup Q)$ is \_\_\_\_\_\_\_\_\_
\begin{qoptions}
\opt{A}{$7$}
\opt{B}{$9$}
\opt{C}{$11$}
\opt{D}{$12$}
\end{qoptions}
\sect{SECTION-II(i)}
\paraq{1}{2}
In $\Delta ABC$, if $a=4$, $b=3$, $c=2$ then
\qn{1} The value of $2\cos B+\cos C$ is
\qn{2} The value of $\dfrac{\sin A+\sin B}{\sin(A+B)}$ is
\paraq{3}{4}
For the equation $\log_{\frac12}|\sin x|=2-\log_{\frac12}|\cos x|$
\qn{3} The sum of solution(s) in $[0,\pi]$ is $k\pi$, then the value of $k$ is \_\_\_\_\_\_\_\_\_
\qn{4} The number of solution(s) in $\left[\dfrac{-\pi}2,\dfrac\pi2\right]$ is
\paraq{5}{6}
If $P=4^{\sin^2x}\cdot16^{\cos^2x}$ ; $x\in R$ \& $y=\dfrac{2+\tan^2(\theta)}{\tan^2(\theta)-3}$ ; where $\theta\ne\dfrac{n\pi}2$ , $n\in I$.
\qn{5} The maximum value of $P$ is \_\_\_\_\_\_\_\_\_
\qn{6} The least positive odd integral value of $y$ is \_\_\_\_\_\_\_
\sect{SECTION-II(ii)}
\qn{7} $\dfrac{\tan\left(\frac{5\pi}6+\frac\theta7\right)\cdot\tan\left(\frac{2\pi}3+\frac\theta3\right)\cdot\cot\left(\frac{10\theta}{21}\right)}{\tan\left(\frac{5\pi}6+\frac\theta7\right)+\tan\left(\frac{2\pi}3+\frac\theta3\right)+\cot\left(\frac{10\theta}{21}\right)}$ is equal to (wherever defined, $\theta\neq n\pi$)
\qn{8} If $K=\cos\dfrac\pi{15}\cos\dfrac{2\pi}{15}\cos\dfrac{3\pi}{15}$ $\cos\dfrac{4\pi}{15}\cos\dfrac{5\pi}{15}$ $\cos\dfrac{6\pi}{15}\cos\dfrac{7\pi}{15}$ then the value of $(1024K)$ is
\qn{9} The number of solutions of the equation $3^{\sin2x+2\cos^2x}+3^{1-\sin2x+2\sin^2x}=28$ in $x\in[0,2\pi]$ is equal to
\papersec{2}
\partsec{1}{PHYSICS}
\sect{SECTION-I}
\qn{1} Find the work done by the force $\vec F=6x^2\,\hat i+12y\,\hat j$ N taking the particle from $(1,1)$ to $(2,2)$.
\begin{qoptions}
\opt{A}{18 J}
\opt{B}{24 J}
\opt{C}{32 J}
\opt{D}{49 J}
\end{qoptions}
\qn{2} A particle of mass 5 kg is moving along $x$-axis and its velocity-time graph is shown in figure. Work done by all forces in 2 to 7 seconds is
\qfig{p2_phy_s1_q02_1}
\begin{qoptions}
\opt{A}{262.5J}
\opt{B}{525J}
\opt{C}{362.5J}
\opt{D}{625J}
\end{qoptions}
\qn{3} Power delivered by a constant force on a particle of mass 2 kg varies with displacement as shown in figure. Particle starts from rest initially. Find the velocity of particle when displacement is 10m.
\qfig{p2_phy_s1_q03_1}
\begin{qoptions}
\opt{A}{$(100)^{1/3}$}
\opt{B}{$(150)^{1/3}$}
\opt{C}{$(200)^{1/3}$}
\opt{D}{$(250)^{1/3}$}
\end{qoptions}
\qn{4} Find work done by frictional force on B by A during displacement.
\qfig{p2_phy_s1_q04_1}
\begin{qoptions}
\opt{A}{100 J}
\opt{B}{0 J}
\opt{C}{$-100$ J}
\opt{D}{50 J}
\end{qoptions}
\qn{5} A block of mass 10 kg is projected with speed of 10 m/s on the surface of plank of mass 10 kg initially at rest. Find work done by kinetic friction on A till frictional force between block and plank stops acting. (Friction coefficient between A and B is 0.1)
\qfig{p2_phy_s1_q05_1}
\begin{qoptions}
\opt{A}{125 J}
\opt{B}{$-250$ J}
\opt{C}{$-375$ J}
\opt{D}{400 J}
\end{qoptions}
\qn{6} A chain of mass $M$ is kept on a hemisphere as shown. Find potential energy of chain.
\qfig{p2_phy_s1_q06_1}
\begin{qoptions}
\opt{A}{$MgR$}
\opt{B}{$Mg(\pi R)$}
\opt{C}{$\dfrac{MgR}{2\pi}$}
\opt{D}{$\dfrac{Mg2R}{\pi}$}
\end{qoptions}
\qn{7} Find work done by force $\vec F=y\hat i+x\hat j$ N displace the particle from point $O$ to point `$C$' along the path $OC$.
\qfig{p2_phy_s1_q07_1}
\begin{qoptions}
\opt{A}{5 J}
\opt{B}{6 J}
\opt{C}{7 J}
\opt{D}{8 J}
\end{qoptions}
\qn{8} Power supplied to a particle of mass 2 kg varies with time as $P=\dfrac{4t^2}3$ Watt. Here `$t$' is time in second. If velocity of particle at $t=0$ is 0 m/s, then the velocity of particle at $t=4$ s will be :
\begin{qoptions}
\opt{A}{$\dfrac{16}3$m/s}
\opt{B}{$\dfrac{25}{16}$m/s}
\opt{C}{$\dfrac{25}9$m/s}
\opt{D}{$\dfrac{16}5$m/s}
\end{qoptions}
\qn{9} A particle is attached at one end of a massless rod whose other end is fixed at $O$ as shown in figure. A particle is given minimum velocity at lower most point to complete vertical circular motion about $O$. Find net force on particle when it is at position $P$. Length of rod is $L$.
\qfig{p2_phy_s1_q09_1}
\begin{qoptions}
\opt{A}{$\dfrac{mg}2\sqrt{39}$}
\opt{B}{$\dfrac{mg}2\sqrt{33}$}
\opt{C}{$\dfrac{mg}2\sqrt{47}$}
\opt{D}{$\dfrac{2mg}3\sqrt{55}$}
\end{qoptions}
\qn{10} If potential energy of particle is given by $U=2x^3+3xy+2y^3$, then find out force acting on it at point $(1,2,1)$.
\begin{qoptions}
\opt{A}{$-12\hat i+27\hat j$}
\opt{B}{$12\hat i+27\hat j$}
\opt{C}{$-12\hat i-27\hat j$}
\opt{D}{$12\hat i-27\hat j$}
\end{qoptions}
\qn{11} Potential energy of a system is given by $U(x)=-2x^2+4x^4$. Which of the following is position of unstable equilibrium ?
\begin{qoptions}
\opt{A}{$x=0$}
\opt{B}{$x=+\dfrac12$}
\opt{C}{$x=-\dfrac12$}
\opt{D}{$x=1$}
\end{qoptions}
\qn{12} The centre of mass of a body :
\begin{qoptions}
\opt{A}{Lies at geometric centre}
\opt{B}{Lies always inside the body.}
\opt{C}{Lies always outside the body.}
\opt{D}{Lies within or outside the body.}
\end{qoptions}
\qn{13} A rider wants to ride a motorcycle on the vertical surface of a large cylindrical wooden wall on commonly known as death well. The radius of well is 20 m and coefficient of friction between the tyres of motorcycle and wall is 0.5. The minimum speed motorcyclist must have in order to prevent slipping should be :
\begin{qoptions}
\opt{A}{30 m/s}
\opt{B}{25 m/s}
\opt{C}{20 m/s}
\opt{D}{15 m/s}
\end{qoptions}
\qn{14} A body is thrown with a velocity of 20 m/s at an angle of $60^\circ$ with horizontal. The radius of curvature of its trajectory at $t=\sqrt3$s after the body began to move is :
\begin{qoptions}
\opt{A}{0 m}
\opt{B}{5 m}
\opt{C}{10 m}
\opt{D}{15 m}
\end{qoptions}
\qn{15} A person rolls a small ball with speed $u$ along the floor from point A. If $x=3R$, determine the required speed $u$, so that the ball returns to A after rolling on the circular surface in the vertical plane from B to C and becoming a projectile at `C' :
\qfig{p2_phy_s1_q15_1}
\begin{qoptions}
\opt{A}{$2\sqrt{gR}$}
\opt{B}{$\sqrt{5gR}$}
\opt{C}{$\sqrt{6gR}$}
\opt{D}{$\dfrac52\sqrt{gR}$}
\end{qoptions}
\qn{16} Find acceleration of lower block w.r.t. upper block, if both blocks were at rest initially.
\qfig{p2_phy_s1_q16_1}
\begin{qoptions}
\opt{A}{1 m/s$^2$}
\opt{B}{0.5 m/s$^2$}
\opt{C}{0.25 m/s$^2$}
\opt{D}{0 m/s$^2$}
\end{qoptions}
\qn{17} System is released at $t=0$. Find acceleration of block of mass $m$ at $t=0$.
\qfig{p2_phy_s1_q17_1}
\begin{qoptions}
\opt{A}{$\dfrac{3g}2$}
\opt{B}{$\dfrac{2g}3$}
\opt{C}{$\dfrac{8g}5$}
\opt{D}{$\dfrac{9g}4$}
\end{qoptions}
\qn{18} A block of mass 4 kg is placed on a rough horizontal surface and coefficient of friction between block and surface is 0.75. Minimum force required to pull the block on horizontal surface is :
\begin{qoptions}
\opt{A}{15 N}
\opt{B}{24 N}
\opt{C}{30 N}
\opt{D}{40 N}
\end{qoptions}
\qn{19} A block of mass 10 kg is placed on rough inclined having angle of inclination $37^\circ$ \& coefficient of friction $\mu=\dfrac12$. This 10 kg block is connected with light string with an another mass $M$ as shown in figure. Initially system is at rest. Calculate maximum and minimum value of $M$ so system does not move.
\qfig{p2_phy_s1_q19_1}
\begin{qoptions}
\opt{A}{14 kg and 2 kg}
\opt{B}{10 kg and 2 kg}
\opt{C}{15 kg and 3 kg}
\opt{D}{9 kg and 1 kg}
\end{qoptions}
\qn{20} The acceleration of a particle moving along $x$-axis varies with velocity $(v)$ as $a=3-2v$. If initially body was at rest, then :
\begin{qoptions}
\opt{A}{Velocity varies with time linearly.}
\opt{B}{Velocity varies with time parabolically.}
\opt{C}{Velocity varies with time exponentially.}
\opt{D}{None of these}
\end{qoptions}
\sect{SECTION-II}
\qn{1} The graph shows velocity v/s $x$-coordinate of a particle moving along $x$-axis. Find out acceleration of particle in m/s$^2$ at $x=4$m.
\qfig{p2_phy_s2_q01_1}
\qn{2} A bike moving with uniform speed of 10 m/s is at a distance of 24m from a car, the car starts from rest and moves with a uniform acceleration of 2 m/s$^2$ away from the bike. If car starts at $t=0$, the time (in sec) at which bike will overtake the car is
\qn{3} A ball is projected at $t=0$ s with velocity 20 m/s at an angle $60^\circ$ with vertical. The radius of curvature of the trajectory at $t=1$ s is $R$. Neglect air resistance and take $g=10$ m/s$^2$. The value of $\dfrac R6$ is
\qn{4} $m_1=4$ kg, $m_2=2$ kg, $\mu=0.4$ \& $g=10$ m/s$^2$. Find $M$ (kg) for which $m_2$ will start sliding over $m_1$. (Ground is frictionless)
\qfig{p2_phy_s2_q04_1}
\qn{5} In the given figure, a force of magnitude $20\,t$ is applied on lower block where $t$ is time in sec. Coefficient of static friction between contact surface is 0.6. For what value of `$t$' upper block begins to slip relative to lower block ?
\qfig{p2_phy_s2_q05_1}
\qn{6} A particle is thrown horizontally with an initial speed of 5 m/s at $t=0$. If the angular velocity of its velocity vector at time $t=2$ is $\dfrac2n$, then find the value of `$n$'. [Take : $g=10$ m/s$^2$ \& neglect air drag]
\qn{7} The angular velocity of a particle is given by $\omega=3t-3t^2+5$ S.I. units. The time when angular acceleration becomes zero is $t$ (in milli second). The value of $t$ is
\qn{8} A particle is released from rest from point A on a smooth track as shown in figure. If the value of acceleration at point B is $\sqrt n\,g$. Find the value of `$n$'.
\qfig{p2_phy_s2_q08_1}
\qn{9} A block slides along a rough and smooth track having length of rough part to be equal to 4m and smooth part having vertical height of 4m as shown in figure. Coefficient of kinetic friction is $\mu=0.5$. The block passes through a point B with a speed of 7 m/s. Find the speed (in m/s) of block when it passes the point `A'.
\qfig{p2_phy_s2_q09_1}
\qn{10} A particle is taken from point $O$ to point $P$ (1, 2) through three different points as shown in figure. Path OCP is parabolic having equation $y=2x^2$. The work done by force $\vec F=xy\hat i+x^2y\hat j$ on particle along path OCP is: ($x$ \& $y$ are in meters \& $F$ is in N) (Round off the answer to nearest integer value)
\qfig{p2_phy_s2_q10_1}
\partsec{2}{CHEMISTRY}
\sect{SECTION-I}
\qn{1} $X$ and $Y$ are two elements which form $\mathrm{X_2Y_3}$ and $\mathrm{X_3Y_4}$. If 0.20 mol of $\mathrm{X_2Y_3}$ weighs 32.0 g and 0.4 mol of $\mathrm{X_3Y_4}$ weighs 92.8 g, the atomic weights of $X$ and $Y$ are respectively
\begin{qoptions}
\opt{A}{16.0 and 56.0}
\opt{B}{8.0 and 28.0}
\opt{C}{56.0 and 16.0}
\opt{D}{28.0 and 8.0}
\end{qoptions}
\qn{2} The vapour density of a mixture containing $\mathrm{NO_2}$ and $\mathrm{N_2O_4}$ is 27.6. The mole fraction of $\mathrm{N_2O_4}$ in the mixture is:
\begin{qoptions}
\opt{A}{0.1}
\opt{B}{0.2}
\opt{C}{0.5}
\opt{D}{0.8}
\end{qoptions}
\qn{3} How many litres of oxygen at 1atm \& 273K will be required to burn completely 2.2 g of propane $(\mathrm{C_3H_8})$
\begin{qoptions}
\opt{A}{11.2 L}
\opt{B}{22.4 L}
\opt{C}{5.6 L}
\opt{D}{44.8 L}
\end{qoptions}
\qn{4} Calculate the total number of atoms in 5.6 L of $\mathrm{SO_2}$ gas at $0^\circ$C and 1 atm. :-
\begin{qoptions}
\opt{A}{$\dfrac{N_A}4$}
\opt{B}{$\dfrac34N_A$}
\opt{C}{$N_A$}
\opt{D}{$\dfrac34$}
\end{qoptions}
\qn{5} If an element have average atomic mass 51.7, find \% age abundance of element having mass number 50 and 52 respectively.
\begin{qoptions}
\opt{A}{10\% and 75\%}
\opt{B}{15\% and 85\%}
\opt{C}{25\% and 95\%}
\opt{D}{35\% and 75\%}
\end{qoptions}
\qn{6} Calculate density of a gaseous mixture which consist of $3.01\times10^{24}$ molecules of $\mathrm{N_2}$ and 32 g of $\mathrm{O_2}$ gas at 3 atm pressure and 860 K temperature (Given : $R=1/12$ atm L/mole K)
\begin{qoptions}
\opt{A}{0.6 g/L}
\opt{B}{1.2 g/L}
\opt{C}{0.3 g/L}
\opt{D}{12 g/L}
\end{qoptions}
\qn{7} The set of species that involve all "p" orbitals of central atom in bonding is.
\begin{qoptions}
\opt{A}{$\mathrm{PH_3}$ , $\mathrm{SO_2}$ , $\mathrm{BeF_2}$ , $\mathrm{NH_3}$}
\opt{B}{$\mathrm{SbH_3}$ , $\mathrm{\overset{+}{C}H_3}$ , $\mathrm{NO_3^-}$ , $\mathrm{N_2}$}
\opt{C}{$\mathrm{HIO_3}$ , $\mathrm{SO_3}$ , $\mathrm{NOF_3}$ , $\mathrm{AsH_3}$}
\opt{D}{$\mathrm{H_2Se}$ , $\mathrm{OF_2}$ , $\mathrm{Be(OH)_2}$ , $\mathrm{NCl_3}$}
\end{qoptions}
\qn{8} Which of the following species can have bond order between any of the two atoms greater than 2.\par
(I) $\mathrm{OCN^{\ominus}}$ ,\par
(II) $\mathrm{H_2C^{\ominus}-C=CH_2}$ ,\par
(III) $\mathrm{HN_3}$
\begin{qoptions}
\opt{A}{I only}
\opt{B}{I, II \& III}
\opt{C}{II \& III}
\opt{D}{I \& III}
\end{qoptions}
\qn{9} Which of the following statement about Mendeleev's table is CORRECT?
\begin{qoptions}
\opt{A}{VIII$^{\mathrm{th}}$ group has 9 elements which are further divided in A \& B subgroups}
\opt{B}{Prediction of properties of Ge were done by leaving a gap below Mn}
\opt{C}{Cu and Cs were placed in same group}
\opt{D}{Correction in Atomic weight were done for B, In, U and Pt}
\end{qoptions}
\qn{10} Which of the following ionic size order is incorrect.
\begin{qoptions}
\opt{A}{$\mathrm{Se^{2-}>I^-}$}
\opt{B}{$\mathrm{S^{2-}>N^{3-}}$}
\opt{C}{$\mathrm{C^{4-}>P^{3-}}$}
\opt{D}{$\mathrm{K^+>Sr^{2+}}$}
\end{qoptions}
\qn{11} Elements P \& Q are from 3$^{\mathrm{rd}}$ \& 4$^{\mathrm{th}}$ period respectively.\par
Atom of element P is paramagnetic in ground state and atom of element Q is diamagnetic in ground state. IF "P" is a p-block element \& Q is from `d'-block and "P" forms a T shape compound with Fluorine. Then
\begin{qoptions}
\opt{A}{Number of electrons having $m=+1$ in Q is 8}
\opt{B}{Number of exchange pairs in $\ell=1$ subshells in P is 10}
\opt{C}{$\mathrm{Q^{2+}}$ is paramagnetic}
\opt{D}{$\mathrm{P_2O_3}$ is basic}
\end{qoptions}
\qn{12} Which of the following set of quantum numbers satisfy last electron entered in La?
\begin{qoptions}
\opt{A}{$n=6$ , $\ell=0$ , $m=0$ , $m_s=+\dfrac12$}
\opt{B}{$n=5$ , $\ell=2$ , $m=2$ , $m_s=-\dfrac12$}
\opt{C}{$n=4$ , $\ell=3$ , $m=0$ , $m_s=+\dfrac12$}
\opt{D}{$n=6$ , $\ell=1$ , $m=+1$ , $m_s=-\dfrac12$}
\end{qoptions}
\qn{13} Which of the following statement is correct ?
\begin{qoptions}
\opt{A}{In $\mathrm{BrF_3}$, Maximum three halogen atoms can lie in same plane.}
\opt{B}{In $\mathrm{SF_2(CH_3)_2}$ molecule, all hydrogen atoms which are involved in s-sp$^3$ overlapping, lie in equatorial plane of $\mathrm{SF_2(CH_3)_2}$.}
\opt{C}{In $\mathrm{OSF_4}$, $\mathrm F_{ax}-\hat S-\mathrm F_{eq}$ bond angle is greater than $120^\circ$.}
\opt{D}{Molecules $\mathrm{IOF_5}$ and $\mathrm{XeO_2F_4}$ have similar shape but have different number of lone pairs in whole molecule.}
\end{qoptions}
\qn{14} Consider following three statement(s) :\par
(i) All atoms are in same plane\par
(ii) All atoms have zero lone pair\par
(iii) All atoms are sp$^2$ hybridized\par
In which of the following molecule at least two statements are true :
\begin{qoptions}
\opt{A}{CO}
\opt{B}{$\mathrm{B_3N_3H_6}$}
\opt{C}{$\mathrm{N_2H_4}$}
\opt{D}{$\mathrm{O=C=O}$}
\end{qoptions}
\qn{15} Match list-I with list-II.\par
{\footnotesize\setlength{\tabcolsep}{2pt}\begin{tabular}{|c|c|c|p{0.3\linewidth}|}
\hline
\multicolumn{2}{|c|}{\textbf{List-I}} & \multicolumn{2}{c|}{\textbf{List-II}}\\ \hline
(P) & $\mathrm{H_3P_3O_9}$ & (1) & S--O--S bond is present\\ \hline
(Q) & $\mathrm{H_2S_2O_7}$ & (2) & Di-basic acid\\ \hline
(R) & $\mathrm{H_2S_4O_6}$ & (3) & P--O--P bond is present\\ \hline
(S) & $\mathrm{H_4P_2O_5}$ & (4) & Central atom (S or P) in maximum oxidation state.\\ \hline
\end{tabular}}\par
\begin{qoptions}
\opt{A}{$P\to1,2;Q\to1,3,4;R\to3;S\to4,2$}
\opt{B}{$P\to1,3;Q\to1,3,4;R\to3;S\to1,3$}
\opt{C}{$P\to4,2;Q\to3,2,4;R\to1;S\to2,3$}
\opt{D}{$P\to3,4;Q\to1,2,4;R\to2;S\to2,3$}
\end{qoptions}
\qn{16} Consider the following :\par
1. $\mathrm{H_5IO_6}$\par
2. $\mathrm{SeOCl_4}$\par
3. $\mathrm{PSCl_3}$\par
4. $\mathrm{SF_6}$\par
Select the correct combination of hybridisation states.
\begin{qoptions}
\opt{A}{1. $sp^3d^2$, 2. $sp^3d$, 3. $sp^3d$, 4. $sp^3d^2$}
\opt{B}{1. $sp^3d$, 2. $sp^3d$, 3. $sp^3d$, 4. $sp^3d^3$}
\opt{C}{1. $sp^3d^2$, 2. $sp^3d$, 3. $sp^3$, 4. $sp^3d^2$}
\opt{D}{1. $sp^3d$, 2. $sp^3d^2$, 3. $sp^3$, 4. $sp^3d^2$}
\end{qoptions}
\qn{17} Consider a molecule `$\mathrm{MX_3}$' which have three $p_\pi$--$d_\pi$ bonds and triangular pyramidal structure Then select the CORRECT statement -\par
[`X' belongs to `p' block, consider no co-ordinate bond is present in molecule]
\begin{qoptions}
\opt{A}{molecule is hypovalent}
\opt{B}{central atom has two lone pairs}
\opt{C}{`X' uses it's `d' orbital for $\pi$ bond formation}
\opt{D}{hybridisation of `M' is $sp^3$}
\end{qoptions}
\qn{18} Second ionsiation potential of Li, Be, B is in the order :
\begin{qoptions}
\opt{A}{Li $>$ Be $>$ B}
\opt{B}{Li $>$ B $>$ Be}
\opt{C}{Be $>$ Li $>$ B}
\opt{D}{Be $>$ B $>$ Li}
\end{qoptions}
\qn{19} Consider the following statements.\par
(P) IE (I) of nitrogen atom is more than IE(I) of oxygen atom.\par
(Q) Electron affinity of oxygen is less than sulphur atom.\par
(R) Electronegativity of an atom has no relation with its ionization enthalpy and electron gain enthalpy\par
Which of the above statements are correct?
\begin{qoptions}
\opt{A}{P, Q}
\opt{B}{P, R}
\opt{C}{Q, R}
\opt{D}{P, Q and R}
\end{qoptions}
\qn{20} The molecule that doesn't exist is:
\begin{qoptions}
\opt{A}{$\mathrm{OF_2}$}
\opt{B}{$\mathrm{NF_3}$}
\opt{C}{$\mathrm{SF_6}$}
\opt{D}{$\mathrm{PH_5}$}
\end{qoptions}
\sect{SECTION-II}
\qn{1} 62.5 gm of a mixture of $\mathrm{CaCO_3}$ and $\mathrm{SiO_2}$ are treated with excess of HCl due to which 1.1 gm of $\mathrm{CO_2}$ is produced. What is the mass \% of $\mathrm{CaCO_3}$ in the mixture?\par
$[\mathrm{CaCO_3+HCl\to CaCl_2+CO_2+H_2O}]$
\qn{2} Ratio of number of protons to neutrons in $3.011\times10^{22}$ ions of $\mathrm{D_3O^+}$ is :
\qn{3} Find number of moles of $\mathrm{Na_3PO_4}$ which contain as many ions as are present in 1368 gm of $\mathrm{Al_2(SO_4)_3}$. (Assuming complete dissociation of salt and no reaction with water)
\qn{4} Calculate percentage change in $M_{avg}$ of the mixture, if 1 mole of $\mathrm{PCl_5}$ undergo 50\% decomposition.\par
$\mathrm{PCl_5\longrightarrow PCl_3+Cl_2}$\par
[Round off to nearest integer]
\qn{5} Number of non planar \& polar species that show $p_\pi$-$d_\pi$ bond among the following.\par
\qfigi{p2_chem_s2_q05_1} , \qfigi{p2_chem_s2_q05_2} , $\mathrm{NOF_3}$ , $\mathrm{POF_3}$ , $\mathrm{XeOF_2}$ , $\mathrm{XeO_2F_2}$ , $\mathrm{SF_6}$ , $\mathrm{HN_3}$ , $\mathrm{B_3N_3H_6}$
\qn{6} If $X$ = maximum number of atoms in a plane in $\mathrm{C(CN)_4}$\par
$Y$ = Number of `d' electrons in Pd\par
then $\dfrac YX$ is \_\_\_\_\_\_.
\qn{7} Number of electrons in $_{41}$Nb that satisfy $n+\ell+m=4$ is \_\_\_\_\_\_.
\qn{8} The sum of oxidation state of Fe in $\mathrm{K_3[Fe(CN)_6]}$ and Cr in $\mathrm{CrO_5}$ is "$X$". Monopositive ion of an element "$Y$" has exactly "$2X$" number of electrons in ultimate shell. If the element belongs to 4$^{\mathrm{th}}$ period then number of correct orders among the following is \_\_\_\_\_\_\_.\par
(i) Zn $>$ Y (Atomic size)\par
(ii) Y $>$ Zn ($\mathrm{IE_1}$)\par
(iii) Y $>$ Ag ($\mathrm{IE_1}$)\par
(iv) Y $>$ Na (number of unpaired electrons)\par
(v) $\mathrm{Y_2O>ZnO}$ (basic nature)
\qn{9} Specie(s) that show significant hydrogen bonding among the following are \_\_\_\_\_\_.\par
$\mathrm{HClO_4}$ , Chloral hydrate, o-chlorophenol,\par
$\mathrm{KHF_2}$ , $\mathrm{HPO_3^{2-}}$ , \qfigi{p2_chem_s2_q09_1}, $\mathrm{SO_2(OH)_2}$
\qn{10} Overlapping(s) that may result in formation of $\pi$-bond is \_\_\_\_\_\_\_.(Internuclear axis is $Z$)\par
$d_{xy}+d_{xy}$ , $d_{yz}+d_{yz}$ , $d_{x^2-y^2}+d_{x^2-y^2}$ , $p_z+p_z$ , $d_{z^2}+d_{z^2}$ , $d_{xz}+d_{xz}$ , $p_x+p_x$ , $p_y+d_{xy}$
\partsec{3}{MATHEMATICS}
\sect{SECTION-I}
\qn{1} If $\cos A=\dfrac34$, then value of $32\sin\dfrac A2\sin\dfrac{5A}2$ is -
\begin{qoptions}
\opt{A}{11}
\opt{B}{$-11$}
\opt{C}{8}
\opt{D}{16}
\end{qoptions}
\qn{2} If $a+b=3-\cos\theta$ and $a-b=\sin\theta$, then maximum value of `$ab$' is -
\begin{qoptions}
\opt{A}{1}
\opt{B}{4}
\opt{C}{8}
\opt{D}{10}
\end{qoptions}
\qn{3} Value of $\left(\tan11^\circ+\tan13^\circ+\tan11^\circ\tan34^\circ+\right.$ $\left.\tan32^\circ+\tan34^\circ+\tan32^\circ\tan13^\circ\right)$ is
\begin{qoptions}
\opt{A}{0}
\opt{B}{1}
\opt{C}{2}
\opt{D}{4}
\end{qoptions}
\qn{4} If $\displaystyle\sum_{m=1}^{9}\sec\left(\theta+(m-1)\dfrac\pi6\right)\sec\left(\theta+\dfrac{m\pi}6\right)=-\dfrac8{\sqrt3}$ ; $\theta\in\left(0,\dfrac\pi2\right)$, then sum of values of $\theta$ is
\begin{qoptions}
\opt{A}{$\dfrac\pi{12}$}
\opt{B}{$\dfrac{2\pi}3$}
\opt{C}{$\dfrac\pi2$}
\opt{D}{$\dfrac{3\pi}4$}
\end{qoptions}
\qn{5} The number of elements in the set $S=\left\{x\in R:2\cos\left(\dfrac{x^2+x}6\right)=4^x+4^{-x}\right\}$ is :
\begin{qoptions}
\opt{A}{1}
\opt{B}{3}
\opt{C}{0}
\opt{D}{infinite}
\end{qoptions}
\qn{6} For $\theta=100^\circ$, the value of $\sin\theta+\cos\theta$ is -
\begin{qoptions}
\opt{A}{less than 0}
\opt{B}{between 0 and 1}
\opt{C}{equal to 0}
\opt{D}{between 1 and $\sqrt2$}
\end{qoptions}
\qn{7} If $4\sin^4x+\cos^4x=1$, then possible values of $x$ is equal to -
\begin{qoptions}
\opt{A}{$n\pi$, $n\in I$}
\opt{B}{$n\pi\pm\sin^{-1}\left(\dfrac25\right)$, $n\in I$}
\opt{C}{$n\pi+\dfrac\pi6$, $n\in I$}
\opt{D}{none of these}
\end{qoptions}
\qn{8} If $2^{\sec^2x}\sqrt{\tan^2y-\tan y+\dfrac12}\le1$, then number of values of $x$ can be, $x\in[0,\ 2\pi]$ -
\begin{qoptions}
\opt{A}{0}
\opt{B}{1}
\opt{C}{3}
\opt{D}{infinite}
\end{qoptions}
\qn{9} Values of $\theta$ satisfying the equation $\left(\sqrt3-1\right)\cos\theta+\left(\sqrt3+1\right)\sin\theta=\left(\sqrt3-1\right)$ is :-
\begin{qoptions}
\opt{A}{$2n\pi+\dfrac{5\pi}6$}
\opt{B}{$2n\pi-\dfrac{5\pi}6$}
\opt{C}{$n\pi\pm\dfrac{5\pi}{12}$}
\opt{D}{$2n\pi\pm\dfrac{5\pi}{12}$}
\end{qoptions}
\qn{10} Two sides `a' and `b' of triangle ABC are given by the roots of the equation $x^2-4x+1=0$ and the included angle between them is $\dfrac\pi3$, then the square of the third side is -
\begin{qoptions}
\opt{A}{$\dfrac14$}
\opt{B}{4}
\opt{C}{11}
\opt{D}{13}
\end{qoptions}
\qn{11} Suppose $a_1,a_2,\ \dots.,a_n,\dots$ be an arithmetic progression of natural numbers. If the ratio of the sum of the first five terms to the sum of first nine terms of the progression is $5:17$ and $110<a_{15}<120$, then the sum of the first ten terms of the progression is equal to -
\begin{qoptions}
\opt{A}{290}
\opt{B}{380}
\opt{C}{460}
\opt{D}{510}
\end{qoptions}
\qn{12} If expression $x+\dfrac1{x^2}\ (x>0)$ attains its minimum value at $x=\alpha$, then $\alpha^3$ is
\begin{qoptions}
\opt{A}{1}
\opt{B}{2}
\opt{C}{3}
\opt{D}{4}
\end{qoptions}
\qn{13} If $2x^{\log_43}+3^{\log_4x}=27$, then sum of possible values of $\sqrt x$ is equal to :-
\begin{qoptions}
\opt{A}{0}
\opt{B}{2}
\opt{C}{4}
\opt{D}{6}
\end{qoptions}
\qn{14} If $\log_5\log_4\log_3\log_2(N)=0$, then number of digits of $N\times2^{19}$ are --\par
\{where $\log_{10}2=0.3010$\}
\begin{qoptions}
\opt{A}{21}
\opt{B}{31}
\opt{C}{20}
\opt{D}{30}
\end{qoptions}
\qn{15} Number of integral values of $m$ for which roots of quadratic equation $(m-1)x^2+3x+(m-3)=0$ lies on either side of origin, is :-
\begin{qoptions}
\opt{A}{1}
\opt{B}{2}
\opt{C}{3}
\opt{D}{4}
\end{qoptions}
\qn{16} Product of all the integral values, which lie in the range of the expression $\dfrac{16x^2-12x+9}{16x^2+12x+9}$ ; $\forall\ x\in R$, is :-
\begin{qoptions}
\opt{A}{2}
\opt{B}{4}
\opt{C}{6}
\opt{D}{8}
\end{qoptions}
\qn{17} For $x\in(0,\pi)$, the equation $\sin x+2\sin2x-\sin3x=3$ has
\begin{qoptions}
\opt{A}{infinitely many solutions}
\opt{B}{three solutions}
\opt{C}{one solution}
\opt{D}{no solution}
\end{qoptions}
\qn{18} If $\alpha$ is root of $\sqrt{x+3}=-2x$ then the value of $|\alpha|$ is
\begin{qoptions}
\opt{A}{1}
\opt{B}{$\dfrac34$}
\opt{C}{$-1$}
\opt{D}{$\dfrac{-3}4$}
\end{qoptions}
\qn{19} The number of distinct solutions of equation $\dfrac54\cos^22x+\cos^4x+\sin^4x+\cos^6x+\sin^6x=2$ in the interval $[0,2\pi]$ is
\begin{qoptions}
\opt{A}{2}
\opt{B}{4}
\opt{C}{6}
\opt{D}{8}
\end{qoptions}
\qn{20} $(100)^{\log_{10}(\ell nx)}=\ell nx-(\ell nx)^2+1$, then $x$ is (where $\ell nx=\ell n_ex$)
\begin{qoptions}
\opt{A}{Irrational number}
\opt{B}{Rational number}
\opt{C}{Even integer}
\opt{D}{odd integer}
\end{qoptions}
\sect{SECTION-II}
\qn{1} In $\Delta ABC$ if $2\cos\dfrac B2\cos\dfrac C2=\dfrac12+\dfrac{(b+c)}a\sin\dfrac A2$ and $A$ is measured in radian then $\dfrac{6A}\pi$ is
\qn{2} Let $S=\{\theta\in(0,2\pi):7\cos^2\theta-3\sin^2\theta-2\cos^22\theta=2\}$. Then, the sum of roots of all the equations $x^2-2\left(\tan^2\theta+\cot^2\theta\right)x+6\sin^2\theta=0\ \theta\in S$, is\_\_\_\_\_\_\_\_.
\qn{3} If the value of $\cos^248^\circ-\sin^212^\circ+\sin132^\circ.\sin12^\circ$ is $t$, then $\dfrac{4t}{\sqrt5}$ is :
\qn{4} Let $S=\left[-\pi,\dfrac\pi2\right)-\left\{-\dfrac\pi2,-\dfrac\pi4,-\dfrac{3\pi}4,\dfrac\pi4\right\}$. Then the number of elements in the set $A=\left\{\theta\in S:\tan\theta\left(1+\sqrt5\tan(2\theta)\right)=\sqrt5-\tan(2\theta)\right\}$ is \_\_\_\_\_
\qn{5} If the sum of solutions of the system of equations $2\sin^2\theta-\cos2\theta=0$ and $2\cos^2\theta+3\sin\theta=0$ in the interval $[0,2\pi]$ is $k\pi$, then $k$ is equal to \_\_\_\_\_\_.
\qn{6} The sides of a triangle are $x^2+x+1$, $2x+1$ and $x^2-1$. Let $\theta$ be the greatest angle of triangle. If $\theta=\dfrac{a\pi}b$, where $a$, $b$ are coprime numbers, then value of $(a+b)$ is
\qn{7} Let $\alpha,\beta\ (\alpha>\beta)$ be the roots of the quadratic equation $x^2-x-4=0$. If $P_n=\alpha^n-\beta^n$, $n\in N$, then $\dfrac{P_{15}P_{16}-P_{14}P_{16}-P_{15}^2+P_{14}P_{15}}{P_{13}P_{14}}$ is equal to \_\_\_\_\_.
\qn{8} If $2\sin\left(\dfrac\pi{22}\right)\sin\left(\dfrac{3\pi}{22}\right)\sin\left(\dfrac{5\pi}{22}\right)\sin\left(\dfrac{7\pi}{22}\right)\sin\left(\dfrac{9\pi}{22}\right)$ is equal to $\lambda$, $\dfrac1\lambda$ is
\qn{9} If $\sin^2x+\sin^2\left(x+\dfrac\pi3\right)+\cos x\cos\left(x+\dfrac\pi3\right)=\dfrac pq$ , $p$ and $q$ are coprime, then $p+q$ is
\qn{10} If $y^{3/8}=\sqrt{\sqrt[3]{3}\cdot\sqrt[3]{9\cdot\sqrt[3]{27\cdot\sqrt[3]{81\cdots}}}}\ \dots\dots\infty$, then $y$ is equal to
\end{multicols}
\end{document}
